\documentclass{article}
\usepackage[T1]{fontenc}
\usepackage{lmodern}
\usepackage{graphicx}
\usepackage{amsmath,amssymb}
\usepackage[a4paper,margin=2cm]{geometry}
\usepackage[absolute,overlay]{textpos}
\setlength{\TPHorizModule}{1mm}
\setlength{\TPVertModule}{1mm}
\usepackage[indonesian]{babel}
\usepackage{hyperref}
\setlength{\emergencystretch}{2em}
% unit: Halaman 22

\begin{document}

% === Halaman 22 (python/native) ===

22

Kategori cipher berbasis bit

1.

Cipher Alir (Stream Cipher)


- beroperasi pada bit (atau byte) secara individual


- enkripsi/dekripsi pesan secara bit per bit (atau byte per byte) hanya dengan

          operasi XOR saja

2.

Cipher Blok (Block Cipher)


- beroperasi pada blok-blok bit (sekumpulan bit)


  (contoh: 128-bit/blok = 16 byte/blok)


- enkripsi/dekripsi pesan secara blok per blok bit dengan komputasi yang lebih

            kompleks dan dilakukan berulang kali dalam sejumlah putaran.

\end{document}
